\chapter{Теоретическая часть. Описание алгоритмов} 
\section{Шум Перлина.}
Шум Перлина был придуман Кеном Перлином в 1982 году. Это функция $Perlin(\overline{x},seed): x \in \mathbb{R}^n \rightarrow z\in [0,1]$. В реальности Шум вычисляется для $n=2, n=3$ из-за большой сложности $o(n  2^n)$

Ход Алгоритма:

1. Пространство $\mathbb{R}^n$ разбивается на n-мерные кубы.

2. Каждой вершине, исходя из seed, присваивается некоторый вектор. 

Для $n=2$ это вектора $$\{\{0,1\},\{0,-1\},\{1,0\},\{-1,0\}\}$$ Для $n=3$ : $$
\left\{
\begin{aligned}
&\{1,1,0\},\{-1,1,0\},\{1,-1,0\},\{-1,-1,0\},\\
&\{1,0,1\},\{-1,0,1\},\{1,0,-1\},\{-1,0,-1\},\\
&\{0,1,1\},\{0,-1,1\},\{0,1,-1\},\{0,-1,-1\}
\end{aligned}
\right\}
$$
Примечание: Можно использовать другой набор векторов. Но у этих наборов есть несколько плюсов:

а) вектора одной длины

b) "равномерно" разбивают пространство


3. Определяется $\overline{x}_{min}=\lfloor \overline{x} \rfloor$ .Это будет вершина n-мерного куба, к которому принадлежит $(x,y)$. Эта вершина однозначна определена для каждого куба (Например, в 2D- левая нижняя вершина)

4. Определяются $y_i, i=\overline{1,2^{n}}$ - вектора в вершинах $z_i$ куба. Также вычисляются $w_i=\overline{x}- z_i$ -вектора из  вершины куба в нашу точку. Затем вычисляется $<y_i,w_i>$.

5.Чтобы Шум не выглядел очень линейным,  используют функцию $\phi : x\in \mathbb{R} \rightarrow y \in \mathbb{R}$, которая должна обладать следующими свойствами:
\begin{enumerate}[label=\alph*)]
    \item $\phi(0)=0;$
    \item $\phi(1)=1;$
    \item $\phi(0.5)=0.5;$
    \item   имеет нулевую производную на конечных точках
\end{enumerate}

В данной реализации используем $\phi(x)=6x^5-1x^4+10x^3$.Теперь будет происходить интерполяция значений $<y_i,w_i>$ поочередна по осям. Например, для 3D будет 8 значений $<y_i,w_i>$, после чего получим 4 значения после интерполяции по оси X. потом эти 4 значения интерполируем по оси Y, получим 2 вектора. Интерполируем их по оси Z.

Получившееся значение и будет значением шума Перлина в нашей точке
\section{Симплексный шум.}
Кен Перлин разработал алгоритм в 2001 году, чтобы устранить ограничения своей классической функции шума, особенно в более высоких измерениях. Это функция $Simplex(\overline{x},seed): x \in \mathbb{R}^n \rightarrow z\in [0,1], \overline{x}=(x_1,x_2,...,x_n)$

Основным преимуществом Симплексного шума над шумом Перлина заключается в том, что вместо разбиения пространства на n-мерные кубы, мы разбиваем его симплексы, что уменьшает количество вычислений с $2^n$ до $n+1$, ($n$).

Ход алгоритма:

1. Вычисляем константы, которые потребуются для преобразований координат: $$F=\frac{\sqrt{n+1}-1}{n}, G=\frac{1-\frac{1}{\sqrt{n+1}}}{n}$$

2.Вычисляем $\overline{x}'=\overline{x}+\sum_{i=1}^nx_i$

Затем вычисляем $ \overline{x}_{loc}= \overline{x}'- \lfloor \overline{x}' \rfloor$

3. Значения внутренней координаты $ \overline{x}_{loc}$ сортируются в порядке убывания, чтобы определить, в каком симплексе м лежит точка.Полученный симплекс состоит из вершин соответствующий упорядоченному обходу ребер от (0, 0, …, 0) до (1, 1, …, 1).Например, точка (0,4, 0,5, 0,3) будет лежать внутри симплекса с вершинами (0, 0, 0), (0, 1, 0), (1, 1, 0), (1, 1, 1).

4. Определяем вектора в вершинах данного симплекса

5. Вычисляем вклад каждой вершины, суммируем и получаем значение в точке. 

\section{Diamond-Square алгоритм.}

Идея впервые была введена Фурнье, Фусселем и Карпентеромна конференции siggraph 1982. Алгоритм начинает работу с 2D сетки, затем, из четырёх начальных значений, случайным образом генерирует карту высот, упорядоченную в виде сетки из точек так, чтобы вся плоскость была покрыта квадратами.

Ход алгоритма:

Алгоритм diamond-square начинает работу с двумерного массива размера 2n + 1. 

1.В четырёх угловых точках массива устанавливаются начальные значения высот. 

2.Шаг diamond. Для каждого квадрата в массиве, устанавливается срединная точка, которой присваивается среднее арифметическое из четырёх угловых точек плюс случайное значение.

3.Шаг square. Берутся средние точки граней тех же квадратов, в которые устанавливается среднее значение от четырёх соседних с ними по осям точек плюс случайное значение.

4.Шаги diamond и square выполняются поочередно до тех пор, пока все значения массива не будут установлены.

\section{Клеточный автомат.}

Клеточный автомаат — дискретная модель, изучаемая в математике. Основой является пространство из прилегающих друг к другу клеток (ячеек), образующих решётку. Каждая клетка может находиться в одном из конечного множества состояний (например, 1 и 0). Решётка может быть любой размерности, бесконечной или конечной (в нашем случае решётка будет конечной и закольцованной), для решётки с конечными размерами часто предусматривается закольцованность при достижении предела (границы). Для каждой клетки определено множество клеток, называемых окрестностью. Мы будем рассматривать 2 вида окрестностей:
\begin{enumerate}
    \item Окрестность фон Неймана, которая задаётся как
    \item Окрестность Мура, которая задаётся как
\end{enumerate}

Смысл использования клеточных автоматов в генерации ландшафтов/пещер - спустя некоторое количество  эпох, применяя определенное правило, получить изображение, похожее на ландшафт/пещеру.

\chapter{Практическая часть}
\section{Модель скрытых переменных.}



